\documentclass[12pt]{article}
\usepackage[margin=1in]{geometry}
\usepackage{amsmath,amsthm,amssymb,scrextend}
\usepackage{array}
\setlength{\parindent}{0pt}

\begin{document}
\newcommand{\ds}{\displaystyle}

\textbf{(i)}

Pardeep twice, Quentin once: 
$$P_{PPQ}=q\left(1-(1-p)^2\right)=q(2p-p^2)$$

Quentin twice, Pardeep once: 
$$P_{PQQ}=p\left(1-(1-q)^2\right)=p(2q-q^2)$$

$$P_{PPQ}-P_{PQQ}=q(2p-p^2)-p(2q-q^2)=pq(q-p)>0$$

$$\therefore \text{Rosalind should play Pardeep twice.}$$

\bigskip\hrule\bigskip

\textbf{(ii)}

Strategy 1:

$$P_1=q\left(1-(1-p)^3\right)=pq(3-3p+p^2)$$

Strategy 2:

$$P_2=\left(1-(1-p)^2\right)\left(1-(1-q)^2\right)=pq(2-p)(2-q)$$

Strategy 3:

$$P_3=p\left(1-(1-q)^3\right)=pq(3-3q+q^2)$$

$$P_1-P_3=pq(3-3p+p^2)-pq(3-3q+q^2)=pq(q-p)(3-p-q)>0$$

\begin{align*}
P_1-P_2
&=pq(3-3p+p^2)-pq(2-p)(2-q)\\
&=pq\left(p-1+(q-p)(2-p)\right)
\end{align*}

For $\ds q-p>\frac12$,

$$p-1+(q-p)(2-p)>p-1+\frac12(2-p)=\frac p2$$

$$\implies P_1-P_2>0$$

$$\therefore \text{Rosalind should choose strategy 1.}$$
For $\ds q-p<\frac12$, let $\ds p=\frac12,\  q=\frac34$

$$p-1+(q-p)(2-p)=\frac12-1+\frac14\cdot\frac32=-\frac18<0$$

$$\implies P_1<P_2$$

Strategy 2 is better than strategy 1

Let $\ds p=\frac12,\  q=\frac78$

$$p-1+(q-p)(2-p)=\frac12-1+\frac38\cdot\frac32=\frac1{16}>0$$

$$\implies P_1>P_2$$

Strategy 2 is worse than strategy 1

\bigskip
\hrule
\begin{center}
    This solution was submitted by \textbf{Jon}\\[4pt]
    Published and refined by \textbf{Jon}\\[8pt]
    \textit{Last updated: July 25, 2026}
\end{center}

\end{document}